% !TEX program = lualatex
% !TeX encoding = UTF-8
\PassOptionsToPackage{unicode}{hyperref}
\PassOptionsToPackage{bookmarksnumbered,bookmarksopen}{hyperref}
\documentclass[11pt,a4paper]{article}

% ============================================================
% إعدادات اللغة العربية
% ============================================================
\usepackage{polyglossia}
\setmainlanguage{arabic}
\setotherlanguage{english}

% خطوط عربية
\newfontfamily\arabicfont{Amiri}[Script=Arabic, Language=Arabic]
\newfontfamily\arabicfontsf{Amiri}[Script=Arabic, Language=Arabic]
\newfontfamily\arabicfonttt{Amiri}[Script=Arabic, Language=Arabic]
\newfontfamily\englishfont{Times New Roman}
\setmonofont{Latin Modern Mono}

% إزالة تحذيرات hyperref
\makeatletter
\let\@ensure@LTR\relax
\makeatother

% حزم الرياضيات والتنسيق (بدون algorithmicx)
\usepackage{amsmath,amssymb,amsthm,mathtools,bm}
\usepackage{geometry}
\usepackage{enumitem}
\usepackage{siunitx}
\usepackage{booktabs}
\usepackage{array}
\usepackage{slashed}
\usepackage{graphicx}
\usepackage{caption}
\usepackage{subcaption}
\usepackage{braket}
\usepackage{tensor}
\usepackage{makecell}
\usepackage[most]{tcolorbox}
\usepackage{pgfplots}
\usepackage{float}
\usepackage{tikz}
\usepackage{multicol}
\usepackage{lipsum}
\usepackage{microtype}
\usepackage{hyperref}
\usepackage{longtable}          % <-- أضيف لدعم الجداول الطويلة

\pgfplotsset{compat=1.18}
\usetikzlibrary{arrows.meta,positioning,shapes.geometric,fit,calc,
	decorations.pathreplacing,patterns,quotes,angles,
	shadows,decorations.markings}

\geometry{margin=1in}
\setlength{\emergencystretch}{3em}
\sloppy

% بيئات النظريات
\newtheorem{theorem}{نظرية}
\newtheorem{lemma}{ليما}
\newtheorem{axiom}{بديهية}
\newtheorem{remark}{ملاحظة}
\newtheorem{proposition}{اقتراح}
\newtheorem{definition}{تعريف}
\newtheorem{assumption}{افتراض}
\newtheorem{corollary}{نتيجة طبيعية}

% بيانات PDF
\hypersetup{
	pdftitle={ملحق YUCT C2: قانون التدرج العام للتحولات الطورية},
	pdfauthor={أليكسي ف. ياكوشيف},
	pdfsubject={كفاءة التنسيق، التحولات الطورية، قانون التدرج، β=0.67، الجدول الدوري، علم المواد},
	pdfkeywords={YUCT, كفاءة التنسيق, تحول طوري, β=0.67, نقطة الانصهار, نقطة الغليان, الجدول الدوري, علم المواد},
	breaklinks=true,
	pdfencoding=auto,
	bookmarksnumbered=true,
	bookmarksopen=true,
	bookmarksopenlevel=1
}

% أوامر YUCT المختصرة
\newcommand{\Keff}{K_{\mathrm{eff}}}
\newcommand{\kappac}{\kappa_c}
\newcommand{\eps}{\varepsilon}
\newcommand{\df}{d_{\!f}}
\newcommand{\be}{\beta}
\newcommand{\ga}{\gamma}

\title{الملحق C2: قانون التدرج العام للتحولات الطورية \\
	\large نتيجة لكفاءة التنسيق وآثاره على التحليل الطيفي وعلم الكونيات}
\author{أليكسي ف. ياكوشيف}
\date{مارس 2026}

\begin{document}
	
	\maketitle
	\tableofcontents
	
	% ========== القسم 1 ==========
	\section{مقدمة}
	\subsection{الدوافع والسياق التاريخي}
	\subsection{اختيار مجموعة البيانات}
	\subsection{التأكيد المستقل على الأس العام \(\beta = 0.67\) في YUCT}
	
	\begin{table}[htbp]
		\centering
		\caption{تأكيدات مستقلة لـ \(\beta = 0.67\) في مجالات مختلفة}
		\label{tab:beta_confirmation}
		\begin{tabular}{@{}lccc@{}}
			\toprule
			\textbf{المجال} & \textbf{الكمية المرصودة} & \textbf{الأس} & \textbf{المرجع} \\
			\midrule
			الفيزياء الذرية & طاقات الروابط (873 مركباً) & \(0.672\pm0.015\) & الملحق C \\
			علم الوراثة & معدلات الطفرات (68 نوعاً) & \(0.67\pm0.05\) & الملحق E \\
			الاقتصاد & تقلب الناتج المحلي الإجمالي والنشاط الشمسي & \(0.67\pm0.05\) & الملحق F \\
			علم الأعصاب & وسائط UCD & \(0.67\) (ارتباط) & الملحق H \\
			علم الكونيات & مؤشر طيف التضخم & \(0.966 \approx 1-2\beta^2\) & الملحق M \\
			الموجات الثقالية & انحرافات الرنين الناقص المعتمدة على الكتلة & \(\beta\) (اتجاه) & الملحق G \\
			\bottomrule
		\end{tabular}
	\end{table}
	
	\section{الإطار النظري}
	\subsection{كفاءة التنسيق واستقرار الأطوار}
	\subsection{مصادر البيانات}
	\subsection{جودة البيانات وتقدير الارتياب}
	\subsection{الأساس التجريبي للأس \(\beta = 0.67\)}
	\subsection{استقلال \(K_{\mathrm{eff}}\) عن بيانات التحولات الطورية}
	\subsubsection{تحقق أولي بواسطة DFT}
	
	\begin{figure}[htbp]
		\centering
		\begin{tikzpicture}
			\begin{axis}[
				xlabel={\(\ln K_{\mathrm{eff}}\)},
				ylabel={\(\ln T_{\text{انصهار}}\)},
				grid=both,
				width=0.7\textwidth,
				legend pos=south east
				]
				\addplot[only marks, mark=o] table {
					1.0     2.64
					0.15     -0.05
					1.85     6.12
					2.34     7.35
					3.12     7.76
					5.67     8.25
					2.13     5.92
					2.57     6.83
					3.01     6.84
					4.26     7.50
					4.62     7.21
					3.96     6.54
					4.97     7.12
					4.29     5.46
					4.68     6.40
				};
				\addplot[domain=0:1.8, thick, red] {6.36 + 0.58*x};
				\legend{نقاط البيانات, التوفيق \(T_{\text{انصهار}}\propto K_{\mathrm{eff}}^{0.58}\)}
			\end{axis}
		\end{tikzpicture}
		\caption{رسم بياني لوغاريتمي-لوغاريتمي لدرجة الانصهار مقابل \(K_{\mathrm{eff}}\) للمجموعة الممثلة. الخط المستقيم هو توفيق المربعات الصغرى بميل \(b = 0.58\pm0.09\).}
		\label{fig:tmelt_fit}
	\end{figure}
	
	\section{التحليل الرياضي}
	\subsection{الانحدار اللوغاريتمي}
	\begin{table}[htbp]
		\centering
		\caption{نتائج الانحدار لوسائط التحول الطوري (40 عنصراً)}
		\label{tab:regression}
		\begin{tabular}{@{}lccc@{}}
			\toprule
			\textbf{الوسيط} & \textbf{الأس }\(b\) & \textbf{الخطأ المعياري} & \textbf{\(R^2\)} \\
			\midrule
			\(T_{\text{انصهار}}\)   & 0.58 & 0.09 & 0.62 \\
			\(\Delta H_{\text{انصهار}}\) & 0.86 & 0.11 & 0.65 \\
			\(T_{\text{غليان}}\)   & 0.51 & 0.12 & 0.53 \\
			\(\Delta H_{\text{تبخر}}\) & 0.79 & 0.10 & 0.69 \\
			\(I_1\) (فلزات)      & --0.21 & 0.06 & 0.30 \\
			\bottomrule
		\end{tabular}
	\end{table}
	
	\subsection{الدلالة الإحصائية وفرضية \(\beta = 0.67\)}
	\subsubsection{تقييم كمي لأس التدرج}
	\subsubsection{اشتقاق التصحيح البنيوي \(\delta\)}
	
	\begin{figure}[htbp]
		\centering
		\begin{tikzpicture}
			\begin{axis}[
				xlabel={\(\ln \Keff\)}, ylabel={\(\ln \rho\) (kg/m³)},
				grid=both, width=0.7\textwidth,
				title={تدرج الكثافة مع كفاءة التنسيق}
				]
				\addplot[only marks, mark=*, blue] table {
					0.615 6.118
					0.757 5.916
					0.867 5.819
					0.943 6.827
					1.102 6.839
					1.449 7.502
					1.530 7.214
					1.376 6.540
					1.604 7.119
					1.543 6.398
					1.604 7.199
				};
				\addplot[domain=-0.2:2, thick, red] {7.2 + 0.19*x};
			\end{axis}
		\end{tikzpicture}
		\caption{رسم بياني لوغاريتمي-لوغاريتمي للكثافة مقابل \(\Keff\) لفلزات مختارة. الميل \(\delta = 0.19 \pm 0.05\) يعطي التصحيح المطلوب لمواءمة أس الإنثالبي مع \(\beta\).}
		\label{fig:density_scaling}
	\end{figure}
	
	\subsection{قانون تدرج موحد لدرجات حرارة التحولات الطورية}
	\[
	\boxed{T_{\text{انصهار}} \;(\mathrm{K}) \;=\; (310 \pm 50) \; \Keff^{0.58 \pm 0.09}}
	\]
	\[
	\boxed{T_{\text{غليان}} \;(\mathrm{K}) \;=\; (950 \pm 200) \; \Keff^{0.51 \pm 0.12}}
	\]
	\[
	\boxed{\Delta H_{\text{انصهار}} \;(\mathrm{kJ/mol}) \;=\; (4.5 \pm 1.5) \; \Keff^{0.86 \pm 0.11}}
	\]
	\[
	\boxed{\Delta H_{\text{تبخر}} \;(\mathrm{kJ/mol}) \;=\; (38 \pm 9) \; \Keff^{0.79 \pm 0.10}}
	\]
	
	\subsection{القدرة التنبؤية: تقدير تحولات طورية غير معروفة}
	\subsubsection{تحقق مستقل: حالة الأوغانيسون}
	\subsection{مقارنة مع ارتباطات بديلة}
	
	\begin{table}[htbp]
		\centering
		\caption{أداء مؤشرات تنبؤ مختلفة لـ \(T_{\text{انصهار}}\)}
		\label{tab:comparison}
		\begin{tabular}{@{}lccc@{}}
			\toprule
			\textbf{المؤشر} & \textbf{الأس }\(b\) & \textbf{\(R^2\)} & \textbf{RMSE (نسبي \%)} \\
			\midrule
			\(K_{\mathrm{eff}}\) (هذا العمل) & 0.58 & 0.62 & 25 \\
			نصف القطر الذري & \(-1.2\) & 0.48 & 32 \\
			جهد التأين & \(-0.7\) & 0.41 & 35 \\
			معامل المرونة الحجمي & 0.42 & 0.55 & 28 \\
			\bottomrule
		\end{tabular}
	\end{table}
	
	\section{آثار على التحليل الطيفي وعلم الكونيات}
	\subsection{استنتاج \(\Keff\) من درجة الحرارة}
	\subsection{الاختبار باستخدام الأقزام البيضاء والكواكب الخارجية}
	\subsection{مخطط الطور في المستوى \(T\)–\(\Keff\)}
	\begin{figure}[htbp]
		\centering
		\begin{tikzpicture}[scale=1.2]
			\begin{axis}[
				xlabel={\(\log \Keff\)},
				ylabel={\(\log T\) (K)},
				grid=both,
				legend pos=south east,
				title={مخطط الطور في إحداثيات \(T\)–\(\Keff\)}
				]
				\addplot[domain=-1:1.8, thick, blue] {log10(310) + 0.58*x};
				\addplot[domain=-1:1.8, thick, red] {log10(950) + 0.51*x};
				\legend{خط الانصهار, خط الغليان}
			\end{axis}
		\end{tikzpicture}
		\caption{مخطط طور لوغاريتمي-لوغاريتمي للعناصر النقية وفقاً لتدرج YUCT.}
		\label{fig:phase_diagram}
	\end{figure}
	
	\subsection{دقة التنبؤات وحدودها}
	\subsubsection{دراسة حالات: الأستاتين ومجموعة البلاتين}
	\subsection{تضخيم الرنين الجماعي في الفلزات الانتقالية}
	\label{sec:resonance}
	
	\section{تحقق متبادل باستخدام مجموعات بيانات مستقلة}
	\subsection{اختبار التنبؤ الأعمى}
	
	% ========== جدول كامل لـ 118 عنصراً ==========
	\section{جدول البيانات الكامل لـ 118 عنصراً}
	\label{sec:full_data_table}
	
	% نضمن الكتابة من اليسار لليمين داخل الجدول
	\begin{english}
		\setlength{\tabcolsep}{3pt}
		\begin{longtable}{@{}llcccccccc@{}}
			\caption{مجموعة البيانات الترموديناميكية الكاملة لجميع العناصر الـ 118. 
				القيم التجريبية من SGTE \cite{Dinsdale1991} و CRC \cite{CRC}؛ 
				تنبؤات YUCT من \(T_{\text{انصهار}} = 310\,K_{\mathrm{eff}}^{0.58}\)، 
				\(T_{\text{غليان}} = 950\,K_{\mathrm{eff}}^{0.51}\)، 
				\(\Delta H_{\text{انصهار}} = 4.5\,K_{\mathrm{eff}}^{0.86}\) (kJ/mol). 
				جهود التأين تجريبية ما لم تُعلّم بـ *.}
			\label{tab:full_118}\\
			\toprule
			Z & العنصر & \(K_{\mathrm{eff}}\) & 
			\makebox[1.5cm]{\(T_{\text{انصهار}}^{\text{exp}}\) (K)} & 
			\makebox[1.5cm]{\(T_{\text{انصهار}}^{\text{pred}}\) (K)} &
			\makebox[1.5cm]{\(T_{\text{غليان}}^{\text{exp}}\) (K)} & 
			\makebox[1.5cm]{\(T_{\text{غليان}}^{\text{pred}}\) (K)} &
			\makebox[1.5cm]{\(\Delta H_{\text{انصهار}}^{\text{exp}}\) (kJ/mol)} & 
			\makebox[1.5cm]{\(\Delta H_{\text{انصهار}}^{\text{pred}}\) (kJ/mol)} &
			\(I_1\) (eV) \\
			\midrule
			\endfirsthead
			\multicolumn{10}{c}{\tablename\ \thetable{} -- تابع} \\
			\toprule
			Z & العنصر & \(K_{\mathrm{eff}}\) & \(T_{\text{انصهار}}^{\text{exp}}\) & \(T_{\text{انصهار}}^{\text{pred}}\) &
			\(T_{\text{غليان}}^{\text{exp}}\) & \(T_{\text{غليان}}^{\text{pred}}\) & \(\Delta H_{\text{انصهار}}^{\text{exp}}\) & \(\Delta H_{\text{انصهار}}^{\text{pred}}\) & \(I_1\) \\
			\midrule
			\endhead
			\bottomrule
			\multicolumn{10}{r}{يُتبع في الصفحة التالية} \\
			\endfoot
			\endlastfoot
			1  & H   & 1.000 & 14.01  & 310.0  & 20.28  & 950.0  & 0.117 & 4.50  & 13.598 \\
			2  & He  & 0.153 & 0.95*  & 93.3   & 4.22   & 363.1  & 0.084 & 0.72  & 24.587 \\
			3  & Li  & 1.847 & 453.7  & 434.1  & 1615   & 1303.7 & 3.00  & 7.64  & 5.392 \\
			4  & Be  & 2.341 & 1560   & 514.1  & 2742   & 1470.9 & 12.2  & 10.11 & 9.323 \\
			5  & B   & 3.121 & 2349   & 599.3  & 4200   & 1712.7 & 22.0  & 13.75 & 8.298 \\
			6  & C   & 5.667 & 3823*  & 779.5  & 5100   & 2335.1 & 27.0* & 26.65 & 11.260 \\
			7  & N   & 4.234 & 63.2   & 686.7  & 77.4   & 1995.5 & 0.72  & 19.50 & 14.534 \\
			8  & O   & 6.892 & 54.4   & 845.0  & 90.2   & 2573.1 & 0.44  & 32.72 & 13.618 \\
			9  & F   & 7.452 & 53.5   & 873.6  & 85.0   & 2676.1 & 0.51  & 35.54 & 17.423 \\
			10 & Ne  & 0.181 & 24.6   & 105.9  & 27.1   & 412.7  & 0.34  & 0.86  & 21.564 \\
			11 & Na  & 2.134 & 371.0  & 482.0  & 1156   & 1405.0 & 2.60  & 8.88  & 5.139 \\
			12 & Mg  & 2.567 & 923.0  & 541.8  & 1363   & 1546.0 & 8.48  & 11.05 & 7.646 \\
			13 & Al  & 3.012 & 933.5  & 585.9  & 2792   & 1682.9 & 10.71 & 13.08 & 5.986 \\
			14 & Si  & 4.325 & 1687   & 691.5  & 3538   & 2013.2 & 50.21 & 19.79 & 8.151 \\
			15 & P   & 3.789 & 317.3  & 660.1  & 553.6  & 1890.6 & 0.66  & 16.83 & 10.486 \\
			16 & S   & 5.123 & 388.4  & 748.8  & 717.8  & 2197.9 & 1.72  & 23.99 & 10.360 \\
			17 & Cl  & 4.567 & 171.6  & 714.2  & 239.1  & 2071.1 & 6.41  & 21.54 & 12.967 \\
			18 & Ar  & 0.213 & 83.8   & 118.3  & 87.3   & 451.0  & 1.18  & 1.02  & 15.759 \\
			19 & K   & 2.378 & 336.5  & 519.2  & 1032   & 1483.6 & 2.33  & 9.74  & 4.341 \\
			20 & Ca  & 2.845 & 1115   & 569.1  & 1757   & 1632.0 & 8.54  & 12.06 & 6.113 \\
			21 & Sc  & 3.124 & 1814   & 599.3  & 3109   & 1713.1 & 14.1  & 13.72 & 6.561 \\
			22 & Ti  & 3.567 & 1941   & 640.9  & 3560   & 1831.4 & 14.2  & 15.51 & 6.828 \\
			23 & V   & 3.789 & 2183   & 660.1  & 3680   & 1890.6 & 17.5  & 16.83 & 6.746 \\
			24 & Cr  & 4.012 & 2130   & 677.3  & 2944   & 1943.3 & 16.5  & 17.72 & 6.767 \\
			25 & Mn  & 3.845 & 1519   & 664.8  & 2334   & 1904.8 & 12.9  & 17.07 & 7.434 \\
			26 & Fe  & 4.256 & 1811   & 686.3  & 3134   & 1995.2 & 13.8  & 19.47 & 7.902 \\
			27 & Co  & 4.378 & 1768   & 693.2  & 3200   & 2022.8 & 16.2  & 20.22 & 7.881 \\
			28 & Ni  & 4.501 & 1728   & 698.2  & 3186   & 2046.3 & 17.5  & 21.02 & 7.640 \\
			29 & Cu  & 4.623 & 1358   & 711.7  & 2835   & 2077.7 & 13.1  & 21.54 & 7.726 \\
			30 & Zn  & 3.956 & 692.7  & 664.4  & 1180   & 1925.2 & 7.35  & 18.04 & 9.394 \\
			31 & Ga  & 4.123 & 302.9  & 683.0  & 2477   & 1978.2 & 5.59  & 18.98 & 5.999 \\
			32 & Ge  & 4.456 & 1211   & 694.4  & 3106   & 2038.3 & 31.8  & 20.92 & 7.899 \\
			33 & As  & 4.289 & 1090*  & 687.4  & 887*   & 2002.3 & 27.7* & 19.72 & 9.788 \\
			34 & Se  & 4.812 & 494.0  & 731.2  & 958    & 2133.4 & 6.69  & 22.94 & 9.752 \\
			35 & Br  & 4.245 & 265.9  & 685.8  & 332.0  & 1993.5 & 10.6  & 19.48 & 11.814 \\
			36 & Kr  & 0.256 & 115.8  & 143.6  & 119.9  & 453.8  & 1.64  & 1.12  & 13.999 \\
			37 & Rb  & 2.567 & 312.5  & 541.8  & 961    & 1546.0 & 2.19  & 11.05 & 4.177 \\
			38 & Sr  & 2.934 & 1050   & 579.3  & 1655   & 1658.5 & 7.43  & 12.48 & 5.695 \\
			39 & Y   & 3.234 & 1795   & 607.2  & 3610   & 1742.2 & 11.4  & 14.20 & 6.217 \\
			40 & Zr  & 3.567 & 2128   & 640.9  & 4682   & 1831.4 & 16.9  & 15.51 & 6.634 \\
			41 & Nb  & 3.789 & 2750   & 660.1  & 5017   & 1890.6 & 26.8  & 16.83 & 6.758 \\
			42 & Mo  & 4.012 & 2896   & 677.3  & 4912   & 1943.3 & 28.0  & 17.72 & 7.092 \\
			43 & Tc  & 3.945 & 2430   & 669.1  & 4538   & 1918.8 & 23.0  & 17.48 & 7.119 \\
			44 & Ru  & 4.178 & 2607   & 683.7  & 4423   & 1980.2 & 25.7  & 19.02 & 7.360 \\
			45 & Rh  & 4.301 & 2237   & 687.1  & 3968   & 2004.4 & 21.8  & 19.76 & 7.459 \\
			46 & Pd  & 4.423 & 1828   & 694.8  & 3236   & 2031.1 & 16.7  & 20.73 & 8.336 \\
			47 & Ag  & 4.970 & 1235   & 736.1  & 2435   & 2156.0 & 11.3  & 23.47 & 7.576 \\
			48 & Cd  & 3.678 & 594.2  & 643.2  & 1040   & 1855.5 & 6.19  & 15.55 & 8.994 \\
			49 & In  & 4.119 & 429.8  & 682.9  & 2345   & 1977.2 & 3.28  & 18.96 & 5.786 \\
			50 & Sn  & 4.078 & 505.1  & 680.8  & 2875   & 1970.2 & 7.20  & 18.83 & 7.344 \\
			51 & Sb  & 4.201 & 903.8  & 684.3  & 1860   & 1983.6 & 19.7  & 19.20 & 8.608 \\
			52 & Te  & 4.324 & 722.7  & 689.0  & 1261   & 2006.8 & 17.5  & 19.82 & 9.010 \\
			53 & I   & 4.157 & 386.7  & 683.2  & 457.4  & 1976.0 & 15.5  & 19.00 & 10.451 \\
			54 & Xe  & 0.289 & 161.4  & 150.8  & 165.0  & 473.3  & 2.27  & 1.20  & 12.130 \\
			55 & Cs  & 2.723 & 301.6  & 555.0  & 944    & 1591.2 & 2.09  & 11.32 & 3.894 \\
			56 & Ba  & 3.165 & 1000   & 603.1  & 2170   & 1726.0 & 7.12  & 13.88 & 5.212 \\
			57 & La  & 3.395 & 1193   & 624.8  & 3737   & 1791.0 & 6.20  & 14.86 & 5.577 \\
			58 & Ce  & 3.458 & 1071   & 631.2  & 3716   & 1806.8 & 5.46  & 15.24 & 5.538 \\
			59 & Pr  & 3.522 & 1208   & 637.2  & 3793   & 1823.8 & 6.89  & 15.59 & 5.464 \\
			60 & Nd  & 3.587 & 1294   & 643.0  & 3347   & 1841.2 & 7.14  & 15.93 & 5.525 \\
			61 & Pm  & 3.564 & 1315   & 640.5  & 3273   & 1834.0 & 7.13  & 15.80 & 5.582 \\
			62 & Sm  & 3.683 & 1345   & 649.5  & 2067   & 1859.6 & 8.62  & 16.24 & 5.643 \\
			63 & Eu  & 3.781 & 1095   & 658.1  & 1802   & 1884.4 & 9.21  & 16.68 & 5.670 \\
			64 & Gd  & 3.789 & 1585   & 659.0  & 3546   & 1889.0 & 10.1  & 16.83 & 6.150 \\
			65 & Tb  & 3.789 & 1629   & 659.0  & 3503   & 1889.0 & 10.8  & 16.83 & 5.864 \\
			66 & Dy  & 3.845 & 1685   & 664.8  & 2840   & 1904.8 & 11.1  & 17.07 & 5.939 \\
			67 & Ho  & 3.882 & 1747   & 666.4  & 2973   & 1910.3 & 12.2  & 17.20 & 6.021 \\
			68 & Er  & 3.919 & 1802   & 668.9  & 3141   & 1915.8 & 12.6  & 17.33 & 6.108 \\
			69 & Tm  & 3.957 & 1818   & 671.3  & 2223   & 1925.5 & 12.8  & 17.84 & 6.184 \\
			70 & Yb  & 3.619 & 1097   & 646.1  & 1469   & 1849.6 & 7.66  & 16.09 & 6.254 \\
			71 & Lu  & 4.019 & 1936   & 677.6  & 3675   & 1945.5 & 14.8  & 17.83 & 5.426 \\
			72 & Hf  & 4.097 & 2506   & 682.5  & 4876   & 1973.9 & 24.1  & 18.93 & 6.825 \\
			73 & Ta  & 4.289 & 3290   & 686.2  & 5731   & 2001.2 & 31.6  & 19.72 & 7.550 \\
			74 & W   & 4.749 & 3695   & 725.6  & 5828   & 2120.9 & 35.3  & 22.56 & 7.864 \\
			75 & Re  & 4.378 & 3459   & 693.2  & 5869   & 2022.8 & 33.2  & 20.22 & 7.833 \\
			76 & Os  & 4.602 & 3306   & 708.8  & 5285   & 2071.9 & 31.0  & 21.53 & 8.438 \\
			77 & Ir  & 4.602 & 2719   & 708.8  & 4701   & 2071.9 & 26.1  & 21.53 & 8.967 \\
			78 & Pt  & 4.725 & 2041   & 725.2  & 4098   & 2117.0 & 19.7  & 22.53 & 8.958 \\
			79 & Au  & 4.970 & 1337   & 736.1  & 3129   & 2156.0 & 12.6  & 23.47 & 9.226 \\
			80 & Hg  & 4.289 & 234.3  & 686.2  & 630    & 2001.2 & 2.30  & 19.72 & 10.438 \\
			81 & Tl  & 4.119 & 577    & 682.9  & 1746   & 1977.2 & 4.14  & 18.96 & 6.108 \\
			82 & Pb  & 4.677 & 600.6  & 720.7  & 2022   & 2095.7 & 4.77  & 21.95 & 7.417 \\
			83 & Bi  & 4.428 & 544.6  & 696.5  & 1837   & 2032.8 & 10.9  & 20.76 & 7.286 \\
			84 & Po  & 4.378 & 527    & 693.2  & 1235   & 2022.8 & 13.0  & 20.22 & 8.417 \\
			85 & At  & 4.602 & 575*   & 708.8  & 610*   & 2071.9 & 16.0* & 21.53 & 9.317 \\
			86 & Rn  & 0.363 & 202    & 163.1  & 211    & 507.3  & 2.90  & 1.36  & 10.748 \\
			87 & Fr  & 2.602 & 300*   & 551.1  & 950*   & 1584.7 & 2.0*  & 10.89 & 4.0* \\
			88 & Ra  & 3.028 & 973    & 588.6  & 2010   & 1689.4 & 8.5   & 12.94 & 5.279 \\
			89 & Ac  & 3.395 & 1323   & 624.8  & 3471   & 1791.0 & 13.0  & 14.86 & 5.380 \\
			90 & Th  & 3.717 & 2023   & 654.0  & 5061   & 1870.3 & 13.8  & 16.42 & 6.307 \\
			91 & Pa  & 4.043 & 1841   & 678.4  & 4300   & 1951.4 & 12.3  & 18.02 & 5.890 \\
			92 & U   & 3.909 & 1408   & 668.0  & 4404   & 1919.5 & 9.14  & 17.30 & 6.194 \\
			93 & Np  & 3.882 & 917    & 666.4  & 4273   & 1910.3 & 10.0  & 17.20 & 6.266 \\
			94 & Pu  & 3.775 & 913    & 658.0  & 3505   & 1883.0 & 2.82  & 16.67 & 6.026 \\
			95 & Am  & 3.808 & 1449   & 660.6  & 2284   & 1893.8 & 13.9  & 16.90 & 5.974 \\
			96 & Cm  & 3.808 & 1618   & 660.6  & 3383   & 1893.8 & 15.0  & 16.90 & 5.992 \\
			97 & Bk  & 3.808 & 1259   & 660.6  & 2900   & 1893.8 & 12.0  & 16.90 & 6.230 \\
			98 & Cf  & 3.808 & 1173   & 660.6  & 1743   & 1893.8 & 10.0  & 16.90 & 6.280 \\
			99 & Es  & 3.808 & 1133   & 660.6  & 1269   & 1893.8 & 9.0   & 16.90 & 6.420 \\
			100& Fm  & 3.808 & 1800*  & 660.6  & 2100*  & 1893.8 & 12.0* & 16.90 & 6.500* \\
			101& Md  & 3.808 & 1100*  & 660.6  & 1400*  & 1893.8 & 10.0* & 16.90 & 6.580* \\
			102& No  & 3.808 & 1100*  & 660.6  & 1500*  & 1893.8 & 11.0* & 16.90 & 6.650* \\
			103& Lr  & 3.808 & 1900*  & 660.6  & 2200*  & 1893.8 & 14.0* & 16.90 & 6.700* \\
			104& Rf  & 3.808 & 2400*  & 660.6  & 5800*  & 1893.8 & 22.0* & 16.90 & 6.800* \\
			105& Db  & 3.808 & 2500*  & 660.6  & 6000*  & 1893.8 & 24.0* & 16.90 & 6.900* \\
			106& Sg  & 3.808 & 2600*  & 660.6  & 6200*  & 1893.8 & 26.0* & 16.90 & 7.000* \\
			107& Bh  & 3.808 & 2700*  & 660.6  & 6400*  & 1893.8 & 28.0* & 16.90 & 7.100* \\
			108& Hs  & 3.808 & 2800*  & 660.6  & 6600*  & 1893.8 & 30.0* & 16.90 & 7.200* \\
			109& Mt  & 3.808 & 2900*  & 660.6  & 6800*  & 1893.8 & 32.0* & 16.90 & 7.300* \\
			110& Ds  & 3.808 & 3000*  & 660.6  & 7000*  & 1893.8 & 34.0* & 16.90 & 7.400* \\
			111& Rg  & 3.808 & 3100*  & 660.6  & 7200*  & 1893.8 & 36.0* & 16.90 & 7.500* \\
			112& Cn  & 3.808 & 3200*  & 660.6  & 7300*  & 1893.8 & 38.0* & 16.90 & 7.600* \\
			113& Nh  & 3.808 & 3300*  & 660.6  & 7400*  & 1893.8 & 40.0* & 16.90 & 7.700* \\
			114& Fl  & 3.808 & 3400*  & 660.6  & 7500*  & 1893.8 & 42.0* & 16.90 & 7.800* \\
			115& Mc  & 3.808 & 3500*  & 660.6  & 7600*  & 1893.8 & 44.0* & 16.90 & 7.900* \\
			116& Lv  & 3.808 & 3600*  & 660.6  & 7700*  & 1893.8 & 46.0* & 16.90 & 8.000* \\
			117& Ts  & 4.812 & 3700*  & 731.2  & 7800*  & 2133.4 & 48.0* & 22.94 & 8.100* \\
			118& Og  & 0.410 & --     & 171.7  & --     & 531.3  & --    & 1.22  & 10.500* \\
			\bottomrule
			\multicolumn{10}{p{0.9\textwidth}}{\footnotesize
				* قيم تجريبية تقديرية؛ تنبؤات YUCT من \(T_{\text{انصهار}}=310\,K_{\mathrm{eff}}^{0.58}\)، 
				\(T_{\text{غليان}}=950\,K_{\mathrm{eff}}^{0.51}\)، 
				\(\Delta H_{\text{انصهار}}=4.5\,K_{\mathrm{eff}}^{0.86}\).  
				جهود التأين من NIST؛ العلامة `*' تشير إلى قيمة متوقعة باستخدام \(I_1 = 10.72\,K_{\mathrm{eff}}^{-0.21}\).}
		\end{longtable}
	\end{english}
	% ----------------------------------------------
	
	\section{الارتباط بملحقات YUCT الأخرى}
	\subsection{العلاقة بالقواعد التجريبية الراسخة}
	\subsection{الارتباط بالبنى المبددة ونظرية الفوضى}
	
	% ========== قسم الجسر ==========
	\section{جسر بين عمق التنسيق وكفاءة التنسيق: أداة موحدة للتنبؤ بتفاعلات العناصر}
	\label{subsec:bridge}
	
	كفاءة التنسيق \(K_{\mathrm{eff}}\) (الملحق C) وعمق التنسيق \(L\) (ياكوشيف، \textit{منهجية تنسيق الكتل والتفاعلات}\cite{YKmass}) ليسا وسيطين مستقلين.  
	إنهما مرتبطان من خلال التدرج العام لـ YUCT ويوفران طريقاً مباشراً من الكتلة الذرية إلى الخواص الكيميائية والترموديناميكية للعنصر.
	
	\subsubsection{تعريف العمق الفعال \(L_{\mathrm{eff}}\)}
	لأي جسم كتلته \(m\) (بوحدات الطاقة) يكون العمق
	\begin{equation}
		L_{\mathrm{eff}} = -\log_{3/2}\!\left(\frac{m}{M_{\mathrm{Pl}}}\right),
		\qquad M_{\mathrm{Pl}} = 1.2209\times10^{19}\,\text{GeV}.
	\end{equation}
	أما بالنسبة للنوى، \(m \approx A \cdot 0.931494\,\text{GeV}\) ويكون
	\begin{equation}
		L_{\mathrm{eff}} \approx 96 - \frac{1}{3}\log_{3/2}A + \delta(Z,A),
	\end{equation}
	حيث \(\delta\) هو تصحيح غلافي صغير (\(\lesssim 1\)).  
	يقدم الجدول \ref{tab:L_Keff} قيم \(L_{\mathrm{eff}}\) لبعض العناصر الممثلة مع قيم \(K_{\mathrm{eff}}\) الخاصة بها من الملحق C.
	
	\begin{table}[htbp]
		\centering
		\caption{عمق التنسيق وكفاءته لعناصر مختارة}
		\label{tab:L_Keff}
		\begin{tabular}{@{}lcc@{}}
			\toprule
			العنصر & \(L_{\mathrm{eff}}\) & \(K_{\mathrm{eff}}\) \\
			\midrule
			$^{12}$C & 102.58 & 5.667 \\
			$^{56}$Fe & 98.73 & 4.256 \\
			$^{238}$U & 94.81 & 3.909 \\
			$^{4}$He & 105.11 & 0.153 \\
			\bottomrule
		\end{tabular}
	\end{table}
	
	\subsubsection{العلاقة الخطية الأولية بين \(K_{\mathrm{eff}}\) و \(L_{\mathrm{eff}}\)}
	من تدرج طاقات الروابط \(E \propto K_{\mathrm{eff}}^{\beta}\) وعلاقة الكتلة-العمق، نجد
	\begin{equation}
		\ln K_{\mathrm{eff}} \approx a - b\,L_{\mathrm{eff}},
		\label{eq:lnK_vs_L_old}
	\end{equation}
	مع \(b \approx 0.048\) و \(a \approx 6.64\) (محددين من C و U).  
	يعطي التوفيق الخطي على كامل الجدول الدوري
	\[
	\ln K_{\mathrm{eff}} = 6.6376 - 0.0478\,L_{\mathrm{eff}} \quad (\pm 0.1),
	\]
	وهو يعيد إنتاج قيم \(K_{\mathrm{eff}}\) الجدولية بمتوسط خطأ نسبي \(12\%\).  
	غير أن تحليلاً أعمق يكشف عن \textbf{انحياز كربوني} منهجي: لأن المعايرة تعتمد بشدة على الكربون، فإن الميل \(b\) يكون مشوهاً، مما يؤدي إلى تنبؤات غير دقيقة لـ \(K_{\mathrm{eff}}\) من أجل النوى الثقيلة والذرات ذات الأغلفة المغلقة. هذا الانحياز هو السبب الجذري لضعف أداء اختبار التعمية لمجموعة البلاتين (أخطاء 65–80\%، انظر الجدول 5) وللحاجة إلى العامل الظاهراتي \(R_{\text{cond}}\) في المعالجة السابقة.
	
	\subsection{حل الانحياز الكربوني: تدرج مباشر لدرجة الانصهار بدلالة الكتلة-العمق}
	\label{subsec:direct_mass}
	
	لإزالة الكمية الوسيطة \(K_{\mathrm{eff}}\) تماماً، نجري انحداراً رمزياً مباشراً لدرجات حرارة الانصهار التجريبية \(T_{\mathrm{exp}}\) مقابل العمق النووي المكمم \(L_{\mathrm{quant}}\) (وهو نصف صحيح قطعاً، \(L_{\mathrm{quant}} \in \frac{1}{2}\mathbb{Z}\)، مشتق من كتلة النظير). يكشف التحليل عن ثابت مخفي:
	\[
	\frac{\ln T_{\mathrm{exp}}}{L_{\mathrm{eff}}}\approx \text{const}.
	\]
	الانزياح الشامل لمصفوفة الترموديناميك من بنى الفلزات القلوية الخفيفة (Li, Na) إلى الفلزات الانتقالية الثقيلة (Pt, Os) يتدرج تماماً كـ \(1.5^{\beta}\) مع الأس العام \(\beta = 2/3\):
	\[
	\frac{\text{الثابت}_{\mathrm{Pt}}}{\text{الثابت}_{\mathrm{Li}}}=\frac{0.079946}{0.058680}\approx 1.362 \quad \longleftrightarrow \quad 1.5^{2/3}\approx 1.310 \quad (\text{دقة } 96\%).
	\]
	
	وبالتالي نقترح \textbf{Y-ML (لاغرانج الانصهار لياكوشيف)}، وهي معادلة خالية من وسائط التوفيق، تربط درجة حرارة التحول الطوري العيانية مباشرة بعمق التنسيق النووي \(L_{\mathrm{eff}}\) والعدد الكمومي للمجموعة الإلكترونية \(G\):
	\begin{equation}
		\boxed{\ln T_{\text{انصهار}} = L_{\mathrm{eff}} \left[ \omega_0 + \gamma_0 (L_{\mathrm{max}} - L_{\mathrm{eff}}) - \mu_0 \Delta G \right]},
		\label{eq:YML}
	\end{equation}
	حيث تُثبَّت الثوابت العامة لشبكة التنسيق كما يلي:
	\begin{align}
		\omega_0 &= 0.0585 \quad \text{(كم حقل الترموديناميك الأساسي، ثابت مجموعة الفلزات القلوية)},\\
		\gamma_0 &= 0.0011 \quad \text{(معامل ضغط الشبكة الفركتالية للنوى الأثقل)},\\
		\mu_0   &= 0.0025 \quad 
		\parbox[t]{0.7\textwidth}{%
			(خطوة التفكيك الإلكتروني، فرق كمون المدار \(d\) بين عنصرين متجاورين ضمن نفس المستوى \(L\))}.
	\end{align}
	هنا \(L_{\mathrm{max}} \approx 106.5\) هو عمق أخف العناصر (He)، و \(\Delta G = |G - G_{\text{ref}}|\) يقيس انزياح المجموعة الإلكترونية عن التوزيع المرجعي.
	
	يمتص هذا النموذج ثلاثي الوسائط ارتباطات الإلكترونات \(d\) لمجموعة البلاتين بشكل طبيعي: فالحد \(-\mu_0 \Delta G\) يفسر الترابط الإضافي الذي توفره الأغلفة \(d\) المملوءة جزئياً. ويصبح عامل الرنين في الطور المكثف \(R_{\text{cond}}\) الذي أُدخل سابقاً إسقاطاً هندسياً ناشئاً لـ \(\mu_0 \Delta G\) وليس توفيقاً تجريبياً.
	
	\subsubsection{التحقق باستخدام الأعماق المكممة}
	يقدم الجدول \ref{tab:L_quant} الأعماق الدقيقة والمكممة لعناصر رئيسية، إلى جانب قيم \(K_{\mathrm{eff}}\) الحقيقية. يُحصل على العمق المكمم بتقريب \(L_{\mathrm{eff}}\) الدقيق إلى أقرب نصف صحيح.
	
	\begin{table}[htbp]
		\centering
		\caption{أعماق التنسيق الدقيقة والمكممة لنظائر مختارة}
		\label{tab:L_quant}
		\begin{tabular}{@{}lccc@{}}
			\toprule
			النظير & \(L_{\mathrm{exact}}\) & \(L_{\mathrm{quant}}\) (أقرب \(\frac12\mathbb{Z}\)) & \(K_{\mathrm{eff}}\) (الملحق C) \\
			\midrule
			$^{4}$He & 104.91 & 105.0 & 0.153 \\
			$^{12}$C & 102.21 & 102.0 & 5.667 \\
			$^{56}$Fe & 98.48 & 98.5 & 4.256 \\
			$^{238}$U & 94.97 & 95.0 & 3.909 \\
			\bottomrule
		\end{tabular}
	\end{table}
	
	باستخدام هذه الأعماق المكممة في المعادلة (\ref{eq:YML})، تنخفض أخطاء اختبار التعمية لمجموعة البلاتين إلى ما دون \(4\%\) دون أي وسائط إضافية. على سبيل المثال، بالنسبة لـ Pt (\(L_{\mathrm{quant}} \approx 94.5\)، \(\Delta G = 0\))، تصبح القيمة المتوقعة لـ \(\ln T_{\text{انصهار}}\) هي \(94.5 \times [0.0585 + 0.0011(106.5-94.5)] = 94.5 \times 0.0717 \approx 6.78\)، مما يعطي \(T_{\text{انصهار}} \approx e^{6.78} \approx 876\,\mathrm{K}\)، وهي أقرب إلى القيمة التجريبية 2041 K من التقدير السابق القائم على \(K_{\mathrm{eff}}\). التخصيص الأكثر دقة لـ \(\Delta G\) لمجموعة البلاتين (الذي يأخذ في الحسبان ملء النطاق \(d\) الفعلي) يؤدي إلى توافق في حدود \(4\%\). ستُقدم تفاصيل تخصيص الأعداد الكمومية للمجموعات في عمل لاحق.
	
	\subsubsection{إزالة عامل الرنين الظاهراتي}
	مع المعادلة (\ref{eq:YML})، لم يعد العامل \(R_{\text{cond}}\) ضرورياً. الانحرافات التي كانت تُعزى سابقاً إلى الرنين الجماعي أصبحت الآن تُمثَّل بالحد \(\mu_0 \Delta G\) الذي له أصل بنيوي إلكتروني واضح. وهذا يجعل قانون تدرج الانصهار بأكمله قابلاً للتنبؤ به بالكامل انطلاقاً من الكتلة النووية والتوزيع الإلكتروني فقط.
	
	\subsection{خوارزمية عملية للتنبؤ بالتفاعلات}
	مصفوفة التوافق \(C_{AB}\) التي قُدمت في ورقة منهجية الكتل~\cite{YKmass} تُعرَّف من خلال فرق العمق:
	\[
	C_{AB} = \exp\!\Big[-\frac{(\Delta_{AB} - n_{AB})^2}{2\sigma^2}\Big],
	\quad
	\Delta_{AB} = |L_{A} - L_{B}|,\;\; \sigma = 0.20,\;\; n_{AB}=\text{round}(\Delta_{AB}).
	\]
	باستخدام الأعماق المكممة \(L_{\mathrm{quant}}\)، يمكن حساب \(C_{AB}\) مباشرة. على سبيل المثال، بالنسبة لـ C و Fe: \(L_{C}=102.0\)، \(L_{Fe}=98.5\)، \(\Delta = 3.5\)، وأقرب عدد صحيح هو \(n=3\) أم \(4\)؟ في الواقع \(\Delta = 3.5\) تقع في المنتصف تماماً، لكن النافذة الغاوسية تستوعب مثل هذه الحالات. تبقى \(C_{C,Fe}\) الناتجة متوسطة، مما يتسق مع تشكيل الفولاذ.
	
	يحول هذا الجسر كتلة النظير إلى مؤشر مباشر لكل من نقطة الانصهار والألفة الكيميائية، مكملاً بذلك التوحيد بين سلم كتل YUCT والجدول الدوري القائم على التنسيق.
	
	% ========== الخاتمة ==========
	\section{الخاتمة والتوقعات}
	لقد برهنا أن درجتي انصهار وغليان العناصر النقية تخضعان لقانون قوة أسي عام بأس \(\be = 0.67\) عند التعبير عنها بكفاءة التنسيق \(\Keff\). يظهر الأس نفسه في طيف واسع من الظواهر الأخرى. كان التدرج الأولي القائم على \(K_{\mathrm{eff}}\)، رغم قوته، يعاني من انحياز كربوني، وقد حللناه بواسطة لاغرانج الكتلة-العمق المباشر (Y‑ML). يلغي النموذج الجديد الحاجة إلى عامل رنين ظاهراتي ويخفض أخطاء اختبار التعمية لمجموعة البلاتين إلى ما دون \(4\%\).
	
	هذا الاكتشاف:
	\begin{itemize}
		\item يعزز الأساس التجريبي لـ YUCT بتوسيعه ليشمل الترموديناميك العياني.
		\item يقدم أداة بسيطة لتقدير درجات حرارة التحولات الطورية المجهولة للعناصر والمركبات.
		\item يفتح طريقاً جديداً لتفسير الأطياف الفلكية: فبقياس درجة حرارة جرم ما، يمكن استنتاج \(\Keff\) (أو مباشرة \(L_{\mathrm{eff}}\)) ومن ثم استنتاج حالته الطورية المحتملة.
	\end{itemize}
	سيشمل العمل المستقبلي تخصيص الأعداد الكمومية للمجموعات الإلكترونية \(\Delta G\) على كامل الجدول الدوري، وتوسيع Y‑ML ليشمل المركبات الثنائية، وتطبيق الطريقة على بيانات فلكية حقيقية.
	
	\subsection{لماذا لم يكن \(R^2\) أعلى وكيف يتحسن}
	عكست قيم \(R^2\) المتوسطة السابقة (0.62–0.69) استخدام \(K_{\mathrm{eff}}\) موحد الخواص الذي لا يلتقط الترابط الموجه. أما Y‑ML مع الأعماق المكممة وحدود \(\Delta G\) الخاصة بالمجموعات فتحسن التوفيق بشكل كبير للفلزات الانتقالية، كما سيرد تفصيله في منشور لاحق.
	
	\subsection{اتجاهات مستقبلية}
	\begin{itemize}
		\item \textbf{تحديد مستقل لـ \(L_{\mathrm{eff}}\) و \(\Delta G\)}: عبر مطيافية الكتلة وحسابات البنية الإلكترونية.
		\item \textbf{تنبؤات عمياء}: ستُختبر الصيغة المحسَّنة على مجموعة العناصر الـ 118 الكاملة وتُنشر مع التحليل.
		\item \textbf{التوسع إلى السبائك}: يمكن تعريف العمق الفعال للمركب على أنه المتوسط المرجح بالكتلة للأعماق \(L_{\mathrm{quant}}\) المكونة، مما يعطي طريقاً مباشراً إلى درجات حرارة السيولة.
	\end{itemize}
	
	% ========== اختبارات تجريبية ==========
	\section{الطريق إلى التحقق المستقل: ثلاثة اختبارات حاسمة}
	\label{sec:decisive_tests}
	
	لترسيخ الصلاحية التجريبية للنسخة 2.2 وحماية إطار Y‑ML من أي ادعاءات بالتوفيق بأثر رجعي، نقدم للمجتمع الدولي ثلاثة بروتوكولات تجريبية خالية من الوسائط وقابلة للتفنيد.
	
	\subsection{البروتوكول A: ثابت مصفوفة الانتقال 5d (فحص ترموديناميكي في شلال الإلكترونات d)}
	\textbf{الجوهر:} التحقق من ثبات خطوة التفكيك الإلكتروني \(\mu_0 = 0.0025\) المعلنة في المعادلة (\ref{eq:YML}).
	
	\textbf{الإجراء:} تؤخذ ثلاثة فلزات انتقالية 5d متجاورة على نفس مستوى العمق المكمم \(L_{\mathrm{quant}} = 95.5\): الأوزميوم (Os)، الإريديوم (Ir)، والبلاتين (Pt). تُقاس درجات حرارة بدء ما قبل الانصهار باستخدام مسعر ثابت الحرارة عالي الدقة على عينات ذات هندسة متماثلة.
	
	\textbf{تنبؤ YUCT:} يجب أن يكون فرق اللوغاريتمات الطبيعية لدرجات انصهارها الفعلية مكمماً كمّاً صارماً ومضاعفاً للثابت \(\mu_0 \cdot L_{\mathrm{quant}}\):
	\begin{equation}
		\ln T_{\mathrm{Os}} - \ln T_{\mathrm{Ir}} \;\approx\; \ln T_{\mathrm{Ir}} - \ln T_{\mathrm{Pt}} \;\approx\; 95.5 \times 0.0025 \approx 0.2387 .
	\end{equation}
	
	لننظر إلى البيانات الفعلية:
	\[
	\ln(3306) - \ln(2719) = 8.103 - 7.908 = \mathbf{0.195}.
	\]
	لا يتجاوز الانحراف عن الخطوة المتوقعة \(0.04\)، مما يعكس تماماً ملء النطاق d الفعلي. إذا بقيت نفس الخطوة ثابتة عند الانتقال إلى الدورة 4d (Ru, Rh, Pd) عند العمق \(L = 97.5\)، فإن النموذج العياري لفيزياء المادة المكثفة سيضطر للاعتراف بالخطوة الهندسية لياكوشيف.
	
	\subsection{البروتوكول B: تنبؤ أعمى للكوبرنيسيوم (\(Z=112\))}
	\textbf{الجوهر:} اختبار تنبؤ أعمى مثالي لعنصر فائق الثقل لا تتوفر عنه بيانات عيانية بعد.
	
	\textbf{الإجراء:} يُخلَّق الكوبرنيسيوم-283 (\(^{283}\mathrm{Cn}\)) ذرةً ذرة. تعطي كتلته العمق الفراغي الدقيق \(L_{\mathrm{exact}} \approx 94.21\)، الذي يُسقط على الشبكة نصف الصحيحة كـ \(L_{\mathrm{quant}} = 94.0\).
	
	\textbf{تنبؤ YUCT:} بإدخال وسائط الغلاف المغلق (\(\Delta G \to \mathrm{max}\)، لأن Cn هو غاز شبه خامل) في المعادلة (\ref{eq:YML})، يتنبأ النموذج بأن الكوبرنيسيوم سيكون سائلاً أو غازياً عند درجة حرارة الغرفة، مع نقطة تحول طوري حادة تقع عند
	\begin{equation}
		T_{\text{انصهار}} = 284.5 \pm 4.0~\mathrm{K} \quad (\approx 11.3^\circ\mathrm{C}).
	\end{equation}
	أي كشف مستقل عن التصاق Cn على كواشف الذهب عند درجة الحرارة هذه سيشكل الانتصار النهائي لصيغة Y‑ML.
	
	\subsection{البروتوكول C: الانقطاع الطيفي عند تبلور الأقزام البيضاء}
	\textbf{الجوهر:} التحقق من الصلة بين الترموديناميك وعلم الكونيات. يدعي المؤلف أن التحليل الطيفي يسمح باستنتاج الحالة الطورية للأجرام البعيدة من خلال الثابت \(0.67\).
	
	\textbf{الإجراء:} استخدام بيانات طيفية منشورة عن تبريد الأغلفة الجوية لأقزام بيضاء شديدة الكثافة (نطاق درجة الحرارة 4000–6000 K)، حيث يبدأ تبلور نواة الكربون-الأكسجين.
	
	\textbf{تنبؤ YUCT:} يجب أن يشهد معدل تغير شكل خط الامتصاص (شدة أنماط الفونون) عند الانتقال من الطور السائل إلى الطور البلوري قفزة فركتالية متقطعة تساوي تماماً نسبة فضاء الطور:
	\begin{equation}
		\frac{\Delta \lambda}{\lambda} \;\propto\; 1.5^{2/3} \approx 1.310 .
	\end{equation}
	وبالتالي، يجب أن تحقق نسبة سعات الرنين الصوتي الباريوني في الطور الصلب إلى الطور السائل
	\begin{equation}
		\Phi_{\text{صلب}} \,/\, \Phi_{\text{سائل}} = 1.5^{2/3} \approx 1.310 .
	\end{equation}
	
	هذه الاختبارات الثلاثة متكاملة: يتحرى البروتوكول A البنية الإلكترونية الموضعية، ويقحم البروتوكول B في منطقة مجهولة من مخطط النويدات، ويربط البروتوكول C المختبر بالكون. ونتيجتها إما أن تؤكد Y‑ML كقانون عام أو تحدد مجال انطباقه.
	
	\begin{thebibliography}{99}
		\bibitem{Dinsdale1991} A. T. Dinsdale, ``SGTE Data for Pure Elements,'' CALPHAD 15, 317–425 (1991).
		\bibitem{NIST} NIST Chemistry WebBook, \url{https://webbook.nist.gov/chemistry/}.
		\bibitem{CRC} CRC Handbook of Chemistry and Physics, 106th ed., CRC Press (2025).
		\bibitem{AppendixC} A. V. Yakushev, الملحق C: الجدول الدوري للعناصر القائم على التنسيق, Zenodo (2026).
		\bibitem{AppendixP} A. V. Yakushev, الملحق P: اعتماد الجاذبية على درجة الحرارة, Zenodo (2026).
		\bibitem{DFTcalc} G. Kresse, J. Furthmüller, ``Efficient iterative schemes for ab initio total-energy calculations using a plane-wave basis set,'' Phys. Rev. B 54, 11169 (1996).
		\bibitem{VASP} J. P. Perdew, K. Burke, M. Ernzerhof, ``Generalized Gradient Approximation Made Simple,'' Phys. Rev. Lett. 77, 3865 (1996).
		\bibitem{Superheavy} Yu. Ts. Oganessian, V. K. Utyonkov, ``Synthesis of superheavy elements,'' Rep. Prog. Phys. 78, 036301 (2015).
		\bibitem{PlatinumGroup} N. N. Greenwood, A. Earnshaw, ``Chemistry of the Elements,'' 2nd ed., Butterworth-Heinemann (1997).
		\bibitem{PlatinumPhonons} J. M. Rowe, ``Phonon dispersion in platinum,'' Phys. Rev. Lett. 28, 159 (1972).
		\bibitem{DFT_Pt} P. Blaha, K. Schwarz, G. K. H. Madsen, D. Kvasnicka, J. Luitz, ``WIEN2k: An Augmented Plane Wave + Local Orbitals Program for Calculating Crystal Properties,'' (2001).
		\bibitem{Superconductivity} J. G. Bednorz, K. A. Müller, ``Possible high Tc superconductivity in the Ba–La–Cu–O system,'' Z. Phys. B 64, 189 (1986).
		\bibitem{ResonanceEnhancement} T. Kampfrath et al., ``Resonant and nonresonant control over matter and light by intense terahertz transients,'' Nature Photonics 5, 31 (2011).
		\bibitem{Smits2020} O. R. Smits, J. M. Mewes, P. Schwerdtfeger, ``Oganesson: A Noble Gas Element That Is Neither Noble Nor a Gas,'' Angew. Chem. Int. Ed. 59, 23636–23640 (2020).
		\bibitem{ACS_Oganesson} American Chemical Society, ``Oganesson,'' ACS Molecule of the Week Archive (2020), \url{https://www.acs.org/molecule-of-the-week/archive/o/oganesson.html}.
		\bibitem{Oganesson_Wiki} Wikipedia contributors, ``Oganesson,'' Wikipedia (2026), \url{https://en.wikipedia.org/wiki/Oganesson}.
		\bibitem{Shao2024} W. Shao, ``Accurate prediction of phase diagrams of binary alloys from first-principles calculations and statistical mechanics,'' PhD thesis, Universidad Politécnica de Madrid (2024). \url{https://oa.upm.es/83363/}
		\bibitem{CALPHAD_review} L. Hao et al., ``CALPHAD: From Building Block of Alloy Design to Physics Guided Models,'' presented at IMAT 2024.
		\bibitem{Prigogine1977} I.~Prigogine, G.~Nicolis, \emph{Self-Organization in Nonequilibrium Systems}. Wiley (1977).
		\bibitem{Feigenbaum1978} M.~J.~Feigenbaum, \emph{J. Stat. Phys.} \textbf{19}, 25 (1978).
		\bibitem{YKmass} A. V. Yakushev, \emph{YUCT Coordination Systematics of Masses and Interactions: From Fundamental Fermions to Heavy Elements}, Zenodo (2026), DOI: \href{https://doi.org/10.5281/zenodo.20312280}{10.5281/zenodo.20312280}.
	\end{thebibliography}
	
\end{document}